\part{Introduction}
Nous cherchons à travers ce projet multidisciplinaire à construire un moteur de simulation de forces centrales avec pour objectif, une implémentation informatique des simulations en Rust avec raylib. 

\bigbreak

Prenons un système à $N$ masses ponctuelles en interaction gravitationnelle, le problème de Cauchy s'écrit, pour toute masse $i\in[\![1,N]\!]$:
\begin{equation}
 \begin{cases}
 \displaystyle\frac{\mathrm d\vec r_i}{\mathrm dt}=\vec v_i,\\[4pt]
 \displaystyle\frac{\mathrm d\vec v_i}{\mathrm dt}
 =G\sum_{j\ne i}m_j\frac{\vec r_j-\vec r_i}{\|\vec r_j-\vec r_i\|^3},\\[4pt]
 \vec r_i(t_0)=\vec r_{i,0},\qquad \vec v_i(t_0)=\vec v_{i,0}.
 \end{cases}
\end{equation}
L'état possède $6N$ composantes dans l'espace!
En effet, chaque corps est décrit par sa position et sa vitesse. Alors, dans l'espace, cela donne $(x_i,y_i,z_i,v_{x,i},v_{y,i},v_{z,i})$, soit six composantes par corps.

Si $f$ est continue et localement lipschitzienne en $y$, le théorème de Cauchy-Lipschitz assure une solution locale unique.

Mais l'existence d'une solution exacte ne veut pas nécessairement dire qu'on dispose d'une formule explicite pour la calculer! $6N$ composantes représente un défi calculatoire de taille que nous allons chercher à approcher de manière mathématique.

\bigskip

L'objectif final est de simuler $N$ corps isolés en interaction. 

Il existe des solutions exactes pour les problèmes jusqu'à deux corps soumis uniquement à leur attraction gravitationnelle. En effet, cela se ramène à un problème central pour leur mouvement relatif. En revanche, dès trois corps, il n'existe plus de formule explicite générale permettant de calculer simplement les trajectoires, même si des solutions particulières existent. 

On utilisera ici l'intégration numérique pour traiter des conditions initiales générales.

\bigskip

Attention : chaque interaction entre deux corps est centrale, mais leur somme sur un corps n'est généralement pas dirigée vers un centre fixe. Les rappels à centre fixe fournissent ainsi des cas de référence avant le passage au système complet.




